\section{Exact renormalization of the spin sources}
\label{sec:sources}

The estimates obtained so far concern densities and partition functions.
To construct fields, we must also compare local observables at different
cutoffs.  We do this by carrying independent sources through the exact
block integration.  The coefficient of the linear spin source determines
the field normalization.  The coupling prescription of
Section~\ref{sec:shooting} remains the prescription at zero source.

Throughout this section a step has side factor $l=L$ or $l=L_*=5L$ and
maps layer $j$ to layer $j-1$.  We use the scales and graph norms of
Section~\ref{sec:rg}; in particular $k=8$, and
\[
 R_j:=H_j^{-1/10}.
\]
The estimates in this section require no information about the
infinite-volume state.  They apply in bulk and on every compatible
period admitted in Section~\ref{sec:free}.

\subsection{The source class and the exact step}

At each site $x$ let $z_x=(z_x^1,z_x^2,z_x^3)\in\C^3$ be an independent
variable.  For a power series $F(q;z)=\sum_\alpha F_\alpha(q)z^\alpha$,
where $\alpha$ ranges over finitely supported multi-indices in the
site and component variables, its coefficient norm at radius $r$ is
\[
 [F]_{k,j;r}:=\sum_\alpha r^{|\alpha|}[F_\alpha]_{k,j}.
\]
The same convention applies to supremum and activity norms.  A
coefficient of a regular function is measured on its graph domain;
a covering activity is measured by its supremum.  We omit $r$ when
$r=1$.  Absolute coefficient sums form a product majorant: products
are bounded by coefficient convolution before any integration or
graph summation.  In particular, identifying source variables does
not increase this norm.

We adjoin to a density in the class of Section~\ref{sec:rg} the
exponent
\[
 \sum_x z_x\cdot q_x+\sum_X 1_XG_X(q;z)
\]
and source-dependent covering activities $k^{\mathrm{src}}_\lambda(q;z)$.
Every $G_X$ and $k^{\mathrm{src}}_\lambda$ has positive source degree.
The graph domain of $G_X$, its load $s_X$, and its mask $1_X$ are those
of a regular error.  Its complete support includes every source site
appearing in it.  Covering labels retain their nonempty inventories of
bad bonds and their complete spin and source dependencies.  All masks
test spin chords only; they impose no condition on a source variable.

The source lists satisfy
\begin{equation}
 \sup_o\sum_{X\text{ anchored at }o}e^{As_X}[G_X]_{k,j}
       \le R_j,
 \qquad
 \norm{k^{\mathrm{src}}}_{j,A}\le w_j.
 \label{eq:orig-27}
\end{equation}
We may separate a regular list by degree or by its multiset of source
indices.  Each such record is anchored at one of those indices; in
degree one this is the unique source site.  Bulk and nonwinding
degree-one regular functions vanish when all their spin arguments
are equal.  No such normalization is imposed on winding records,
and no alignment normalization is imposed in higher source degree.
In particular, a spin-independent term of degree at least two stays
in its source-supported regular record; it is not extracted as a
volume scalar.
The classifications as short or winding apply to the source indices
and anchors as well as to the spin dependencies.

The lists inherit the spatial covariance conventions of the
source-free class.  In particular they are translation covariant
as lists in independent source variables: translating a complete
record, its anchor, and all of its spin and source indices leaves
its coefficient function unchanged after the corresponding
relabeling of the arguments.  Thus a specialization to finitely
many nonzero sources need not be translation invariant, but the
bulk rule that produces its coefficients is translation covariant.
The preparations, record assignments, and normalization rule are
also reflection covariant as a family of prescriptions.  A
reflection preserves the standard prescription and exchanges
the two mirror inclined prescriptions, with their reflected
input lists.  The same conventions apply to regular and
covering source records.

The lists have real coefficients and are covariant under simultaneous
internal $O(3)$ transformations of spins and sources.  They are
organized by retaining the full vector or tensor polynomial associated
with a multiset of source sites, including all its component
coefficients.  Thus covariance is imposed on a tensor contribution,
not on each scalar component separately.  No tensor norm with an
uncontrolled dependence on the source degree is used.  At the initial
layer the source factor is simply $\exp(\sum_x z_x\cdot q_x)$, so the
two additional lists are zero and all these conditions hold.

\begin{proposition}[Exact source step]
\label{prop:source-step}
Consider an admitted source-free step of Section~\ref{sec:rg} from
layer $j$ to layer $j-1$, with precise coupling
$b\in[H_j/2,2H_j]$, together with source lists satisfying
Equation~\eqref{eq:orig-27} and the normalization and support
conditions above.  There is a real number $c_j>0$, determined by the
bulk linear response, such that
\begin{equation}
 z_x=\frac{z'_{[x]}}{l^2c_j},
 \qquad |c_j-1|\le C_LR_j,
 \label{eq:orig-28}
\end{equation}
gives an exact output representation with leading exponent
$\sum_Yz'_Y\cdot V_Y$ and renewed source bounds
Equation~\eqref{eq:orig-27} at layer $j-1$.  The source-free output
is exactly the output at zero source.  No further source-dependent
volume scalar is extracted.

For two compatible inputs to a regular step, let $u,\lambda,
\epsilon_h$ be the source-free discrepancies of
Equation~\eqref{eq:orig-5}.  Let $u_s$ be a fixed positively weighted
sum of the regular and covering source discrepancies divided by
$R_j$ and $w_j$, respectively; the regular discrepancy is measured
through order $k-1$, with absent records padded by zero.  The weights
can be chosen so that, for some $q_s<1$,
\begin{equation}
 \begin{split}
 u_s'&\le q_su_s+
   C_L(\log H_j)^D
          \bigl(u+\epsilon_h+|\lambda|/H_j\bigr),\\
 |c_{2,j}-c_{1,j}|&\le C_LR_j
    \left[u_s+(\log H_j)^D
          \bigl(u+\epsilon_h+|\lambda|/H_j\bigr)\right].
 \end{split}
 \label{eq:orig-29}
\end{equation}
These estimates hold in bulk and simultaneously on the compatible
admitted periods.  The fixed exponent $D$ may depend on $P_0$.
The order of choices is a sufficiently large $L$, then $P_0$ large
enough for the fixed polynomial losses, and finally sufficiently
large $H$.
\end{proposition}

\begin{proof}
We first integrate after the substitution $z_x=z'_{[x]}/l^2$, leaving
the final scalar division for the end of the proof.  We estimate the
output at source radius $2$; any fixed finite collection of such
radii can be used by increasing $L$.  A coefficient of degree $a$
then gains $(2/l^2)^a$.  Since each block has $l^2$ sites and
$T_x=V_{[x]}$, subtracting the coarse leading exponent gives the
identity
\[
 \frac1{l^2}\sum_x z'_{[x]}\cdot q_x
 =\sum_Yz'_Y\cdot V_Y+
       \frac1{l^2}\sum_xz'_{[x]}\cdot(q_x-T_x).
\]
It remains to integrate the second term together with the old
source lists.

\paragraph{Preparation of the source factors.}
Write $M=\mathfrak m_j$, $p=p_j$, and $g=b^{-1/2}$.  We use the
sector integration of Section~\ref{sec:rg}.  A core is one of its
connected exceptional regions, with load $s_c$ and the Gaussian
read set appearing in Equation~\eqref{eq:orig-7}.  A prepared
letter is the exponential-minus-one factor used in that integration.
There are three additional sorts of factor.

First consider the term containing $q_x-T_x$.  All such factors
whose sites lie within distance $12M$ of a core's seeds are included
in that core factor.  If these neighborhoods overlap, each site is
assigned once, using the ownership convention of
Section~\ref{sec:rg}.  The attached profiles are extended to any
new exterior variables read by these factors.  The absolute source
series contributes at most $\exp(C_LM^Ds_c)$ to the core envelope;
without spin derivatives the bound is
$\exp(C_LM^2s_c)$.  This is paid by the exceptional reserve
$\exp(-c_Lp^2s_c)$ in Equation~\eqref{eq:orig-7}, after the prescribed
choice of $P_0$ and $H$.  These core factors may depend on the actual
spin variables.

Away from those neighborhoods, choose a smooth angular plateau
containing every full-sector value of $q_x-T_x$ and write its
exponential as one plus a letter.  The letter vanishes at angle
zero, and the plateau gives $|q_x-T_x|\le C_Lt_j$.  Consequently
its majorant, through the required normalized spin derivatives, is
\begin{equation}
 C_Lg\operatorname{poly}(M,p).
 \label{eq:orig-30}
\end{equation}
The mean estimates in Section~\ref{sec:rg} give the same order for
those derivatives.  Gaussian polynomials produced by a fixed number
of differentiations are covered by its joint estimates.  The new
reads are the site, its block reference, and the attached profiles;
in the remote preparation they also include the stencil deciding
the selection status.  These are already contained in the halos,
core footprints, and status conventions of the sector construction.
We use its local extension for a term with no output bad-bond flag.

Second, prepare a masked $G_X$ as an old regular error, with the
graph cutoff in the remote case and with exponential-minus-one
expansion.  The preparations in \Rref{rg:prepared} use its small
norm, its graph domain, and its complete supports; they do not use
the affine-Hessian normalization of a source-free error.  In a
linear use, the coefficient majorant retains the substitution
factor $(2/l^2)^a$.  Nonlinear powers are bounded by products.
The summed majorants at a support hit, including each of the fixed
larger output load exponents required below, are at most
$C_L\operatorname{poly}(M,p)R_j$.  A nonremote measurable evaluation
still meets its core.  Changing the anchor of a spin direction to
the chosen source anchor costs only fixed powers of $M$ and the
total load, also in a discrepancy estimate.  No source integration
domain is being translated.

Third, the positive-degree parts of old covering labels are
prepared by the same activity estimates.  The substitution and
Equation~\eqref{eq:orig-27} satisfy those activity hypotheses with
at most a fixed combined factor.  Every such use still meets a
primary exceptional object.  Terms with a nonempty output
inventory require positive-degree supremum bounds only, without
spin derivatives.

\paragraph{Absolute coefficient sums and exact reassembly.}
For clarity, all product estimates are applied after taking
absolute coefficient majorants and performing coefficient
convolution.  The majorant of an exponential is the exponential
series of the majorant of its exponent.  For nonlinear powers of
a masked $G_X$, its undifferentiated envelope is small before
cutoff derivatives are charged.  At most $k$ factors can receive
derivatives.  Their Leibniz placements cost a polynomial in the
number of factors, so the one-factor bounds above remain valid
with an additional
$C_L\operatorname{poly}(M,p,1+s_X)$ multiplier where required.
The passage from anchor sums to support-hit sums uses the fixed
projection powers and spare support exponent exactly as for old
errors.

Repeated source variables cause no change in this argument.
Factorization of separate contributions is determined by their
complete supports, not by disjointness of their source monomials.
Inside a core, coefficient supremum sums of exponent powers have
the core envelope already obtained; differentiation through the
transport produces only its fixed polynomial score costs.
Thus the hypotheses of \Rref{rg:joint-product} and
\Rref{rg:tree-condition} hold in the coefficient norm.  Each
argument mark is retained in the output graph.  A singleton
transfer inherits its source anchor.  Any other connected term is
anchored at one of its source marks.  If covariance requires
several anchor choices, split the term with nonnegative
coefficients summing to one and identical complete records, or
retain its ordered marks.  This operation incurs no loss growing
with source degree.

Apply the exact connected reassembly \Rref{rg:connected} to these
lists.  Keep every zero-degree prescription fixed.  In particular,
the scalar extraction, canonical normalization, and error reset
are first performed at degree zero, including their exact gas
recombinations.  Their output is the source-free output.  The
remaining exact algebra produces regular and covering lists of
positive source degree.  Connected source terms with an output
bad-bond flag retain the stronger activity bound $o(w_{j-1})$.

We next separate the part that needs a contraction estimate from
the terms that are small by order.  A regular output term using
Equation~\eqref{eq:orig-30} is bounded by
$C_Lg\operatorname{poly}(M,p)$.  A term containing a new core
decoration or an old positive-degree covering activity has
arbitrary-power accuracy in $g$.  A term using two or more
$G_X$ factors is bounded by
$C_L\operatorname{poly}(M,p)R_j^2$.  The same bound applies when
one $G_X$ is accompanied by an ordinary source-free correction
or failure factor, since its order in $g$ is smaller than $R_j$
at large $H_j$.

Here these order statements concern summed lists, not just an
individual term.  Outside cores and exceptional terms, the
integrated and Mayer-log expansions have summed source-free
envelopes $O_L(g\operatorname{poly}(M,p))$, together with errors
of still smaller cap.  Inflate each marked majorant by the
inverse of its claimed order in $g$ or $R_j$, divided by
$C_LM^{D'}p^{D'}$ with a sufficiently large fixed $D'$.  The
summed hit bound, including reserved output exponents, remains
$o_L(M^{-2})$.  The tree estimate therefore preserves the marked
powers.  This also treats quadratic and higher terms from the
exponential of a single regular source function.  In a compulsory
core the source series is charged inside the core factor, where
its exceptional reserve remains available; it is not treated as
a small optional envelope.  The same estimate applies to the
subsequent list operations and mask regroupings, using their
spare support exponents.  Those regroupings can carry source
degree only in covering terms, apart from the linear-source
normalization performed below.

The terms not covered by these small-order estimates are the
isolated linear transfers of old $G_X$.  They are the empty-pattern
Gaussian expectations of \Rref{rg:error-transfer}, with the
projected spin formula of Section~\ref{sec:rg}, with $G_X$ in
place of the old error, and with each source index specialized
to its coarse index.  Their old cutoffs, required output graph,
and output mask are retained.

\paragraph{The isolated transfer in higher source degree.}
For degree at least two, the anchored norm of the isolated
transfer, with any fixed required larger output exponent, is
\begin{equation}
 \bigl(C/L^2+o_H(1)\bigr)R_j.
 \label{eq:orig-31}
\end{equation}
There is no Taylor normalization in this assertion.  Each coarse
source anchor has $l^2$ possible old anchors, whereas a degree-$a$
coefficient gains $(2/l^2)^a$.  For $a\ge2$ their product is at
most $4/l^2$.  It remains to check that composition with the spin
map has a leading bound fixed before $L$ is chosen.

Use the Gaussian guard from the proof of
Equation~\eqref{eq:orig-9}, including the normal component.  For
$s_X\le L^{1/4}$, the old cutoff is on its plateau.  The absolute
first jet of $q_x$, measured in the old normalized direction norm,
is bounded by the row of $P_h$ with its weighted moments.  Its
leading constant is independent of $L$.  Differentiating the
tangent projection of the fluctuation gives an extra $t_j$;
higher derivatives of the spin map also have an extra $t_j$,
with fixed powers of $M$ at fixed $L$.  The value itself requires
only a supremum estimate.  The projected short load, including
the necessary output halo, is bounded by a constant independent
of $L$, as in \Rref{rg:error-contraction}.  Thus composition and
the fixed output load weight contribute a fixed multiplier.

For $s_X>L^{1/4}$, the direct composition estimate from
Equation~\eqref{eq:orig-9} has a fixed polynomial in $1+s_X$.
The unused old support exponential pays this polynomial and the
anchor count.  At every load the guard complement costs an
exponentially small Gaussian tail times fixed polynomial losses.
After listing these losses, choose $P_0$ and then $H$ so that they
are negligible.  Summing the old coefficients proves
Equation~\eqref{eq:orig-31}.

\paragraph{The isolated transfer in degree one.}
For degree one the corresponding bound is
\begin{equation}
 \bigl(C/L+o_H(1)\bigr)R_j.
 \label{eq:orig-32}
\end{equation}
High loads are again paid by their spare support exponential.
For a short label, write its degree-one coefficient as a vector
function relative to its old anchor spin.  This vector vanishes
at alignment.  On the same guard, the predicted relative spin
values without Gaussian perturbation have old normalized size
$C/L+o_H(1)$, by the relative-map estimate used in
Equation~\eqref{eq:orig-9} and \Rref{rg:affine-prediction}.
The retained Gaussian perturbation adds at most
$C_L\max|Z|/p$, with fixed powers of $M$ independent of $P_0$.
Its expectation is $o_H(1)$: there are $O_L(M^D)$ read
coordinates, so either a union estimate or the sum of their
fixed marginal moments suffices after the choice of $P_0$.

Apply a mean-value estimate between the aligned configuration
and the guarded spin configuration.  Both the guarded chords
and this interpolation remain strictly in the graph domain
through $4t_j$.  The value of the integrated vector function
therefore has the factor $C/L+o_H(1)$.  This estimate precedes the
anchor summation; the degree-one substitution contributes
$2/l^2$, canceling that summation's $l^2$ factor up to a constant.

For positive-order derivatives, each first derivative of the
relative map has the same $C/L+o_H(1)$ bound on the guard, and
every higher derivative is $o_H(1)$.  Moving the old anchor may
also rotate the vector as a whole.  Internal covariance expresses
this by a local smooth rotation of the resulting vector; every
normalized derivative falling on that rotation contains a
factor $t_j$ times fixed polynomial losses, for which the input
supremum is sufficient.  Thus no derivative above order $k$ is
needed.  This proves Equation~\eqref{eq:orig-32}.

The same two singleton arguments apply to input differences
through order $k-1$.  A change of Gaussian map or kinetic history
acting on an unchanged input uses one additional input derivative,
at most order $k$, and costs
$C_L\operatorname{poly}(M,p)R_j$ times
$u+\epsilon_h+|\lambda|/H_j$.  The estimates therefore retain
their input contraction and the required map forcing.

\paragraph{Difference estimates for the other terms.}
The one-difference joint estimates of Section~\ref{sec:rg} apply
to the prepared source letters as well.  They preserve the order
in Equation~\eqref{eq:orig-30}, with an additional
$C_L\operatorname{poly}(M,p)$ factor multiplying
$u+\epsilon_h+|\lambda|/H_j$.  In particular, changing $g$ costs
$O_L(g|\lambda|/H_j)$ in that angular factor.  For a masked
source function, the same assertion uses its old graph norm
and the composition estimate just described.  Powers of $g^{-1}$
from hard transports are charged only against exceptional
reserves.  Consequently exceptional terms retain arbitrary-power
accuracy or their stronger output-inventory price.

A difference in one of two or more regular source factors gives
the nonlinear bound
$C_L\operatorname{poly}(M,p)R_j^2u_s$.  Marking the relevant
factor in the inflated-majorant argument above proves the same
bound for the entire connected list.  At most $k-1$ derivatives
fall on a changed input, and at most $k$ on an unchanged input
when its map changes.  These are precisely the differentiated
joint and tree hypotheses; no higher derivative is invoked.

\paragraph{Extraction of the linear coefficient.}
Each bulk or corresponding short degree-one output term is
$z'_Y\cdot F_X(V)$, anchored at its source site $Y$.  If all spin
arguments equal $v\in S^2$, covariance gives
$F_X(v,\ldots,v)=a_Xv$ for a real scalar $a_X$.  Indeed the
stabilizer of $v$ fixes only its span, and transitivity makes the
scalar independent of $v$.  Use the exact identity
\[
 1_Xz'_Y\cdot F_X
 =a_Xz'_Y\cdot V_Y
  +1_Xz'_Y\cdot(F_X-a_XV_Y)
  -(1-1_X)a_Xz'_Y\cdot V_Y.
\]
The middle term vanishes at alignment.  Its norm is at most a
fixed multiple of the original norm, since the anchored
derivatives of $V_Y$ are bounded.  The final term is an off-mask
correction and is put in the covering gas by the exact
normalization operation of \Rref{rg:normalization}.

Sum the extracted coefficients per bulk anchor.  Spatial
translation covariance of the source lists makes this sum
independent of the anchor.  Together with
the previous coefficient $1$ of the leading exponent, this gives
$c_j$.  Use that same bulk sum on every period.  Actual winding
terms are left unnormalized.  When a bulk coefficient has been
extracted without a corresponding short term on the period,
retain its compensating one-site term as a winding entry.  A
one-site graph is allowed; alternatively it may be enlarged by
nearby bonds with the corresponding off-mask correction.  A
missing bulk coefficient has complete length at least the
winding threshold, so its spare reserved exponent pays this
declared record.  The classification includes the changes to
the records themselves, exactly as in the source-free
normalization.  Winding spin functions are not required to
vanish at alignment.

Off-mask source corrections have the stated summed bounds with
their complete support prices; passing from an anchor to a
support hit costs only fixed powers of $M$.  Contributions
carrying an output bad-bond seed retain their $o(w_{j-1})$
activity under opening and regrouping.  At zero source every
operation in this paragraph is the identity, so it changes no
source-free label.  Equations~\eqref{eq:orig-30}--\eqref{eq:orig-32}
and their difference versions give $|c_j-1|\le C_LR_j$ and the
second estimate of Equation~\eqref{eq:orig-29}.  In particular
$c_j>0$ for sufficiently large $H$.

Finally divide all output source variables by $c_j$.  The radius-$2$
bound controls this substitution at radius $1$, its differences,
and the sum over all positive degrees.  For a common scalar
dilation, the extra degree factor in a difference is paid by the
spare analytic radius.  The leading exponent is now
$\sum_Yz'_Y\cdot V_Y$.

To close the caps and discrepancies, divide the output bounds
by $R_{j-1}$ and $w_{j-1}$.  The isolated regular multipliers,
including alignment subtraction and the fixed normalization
costs, can be made strictly smaller than one by choosing $L$.
Moreover,
\[
 \frac{R_j}{R_{j-1}}\le1,
 \qquad
 \frac{g}{R_{j-1}}\operatorname{poly}(M,p)\longrightarrow0,
 \qquad
 \frac{R_j^2}{R_{j-1}}\operatorname{poly}(M,p)\longrightarrow0
 \quad(H_j\to\infty).
\]
The normalized gas discrepancy has its $o(w_{j-1})$ allowance.
Changing the dilation on an already small positive-degree
remainder gains the additional factor $R_j$ from the estimate
on $c_j$.  These observations give the first estimate of
Equation~\eqref{eq:orig-29} and renew
Equation~\eqref{eq:orig-27}.  All fixed caps and constants can
be enlarged in the stipulated order; their finite sizes do not
alter the leading constants in
Equations~\eqref{eq:orig-31} and~\eqref{eq:orig-32}.
\end{proof}

\subsection{Endpoint source convergence and the normalization product}

The source step determines the field at every cutoff.  Two
consequences will be used separately: convergence at a fixed
bottom layer, and a uniform comparison of the normalizations
for the ordinary and inclined first steps.

\begin{corollary}[Sources at a fixed bottom layer]
\label{cor:source-endpoints}
For the trajectory comparisons in Equation~\eqref{eq:orig-14},
the source discrepancy $u_s$ tends to zero at each fixed bottom
layer, and the difference of the step coefficients $c_j$ tends
to zero at each fixed bottom step.  In every fixed compatible
finite period, the resulting source exponents and covering
activities converge in their coefficient norms, with the step
dilations above that layer included.
\end{corollary}

\begin{proof}
At the top of a comparison, Equation~\eqref{eq:orig-27} bounds
the normalized source discrepancy by a fixed constant.
Iterating Equation~\eqref{eq:orig-29} gives a geometric factor
$q_s$ on this initial discrepancy and a geometric convolution
of the forcing from Equation~\eqref{eq:orig-14}.  Fixed powers
of the scale indices do not prevent that forcing from tending
exponentially to zero at any fixed bottom layer.  The second
line of Equation~\eqref{eq:orig-29} gives the assertion for
each fixed step coefficient.  On a fixed finite period,
covering expansions converge absolutely, while masked regular
functions are controlled by their supremum bounds.  Thus their
integrated source functions have the asserted comparison.
No claim of uniformity in a growing volume is needed here.
\end{proof}

Write $c_{j,N}$ for the coefficient of the step from layer $j$
in a standard run of depth $N$, and define
\begin{equation}
 B_N:=\prod_{j=1}^Nc_{j,N}^{-1}.
 \label{eq:field-normalization}
\end{equation}
This is the field normalization; the area factor $a_N^2$ will
be included separately when fields are smeared.  The product
is positive and finite.  Since $H_j$ grows linearly in $j$,
Equation~\eqref{eq:orig-28} gives
$|\log B_N|=O_{L,H}(1+N^{9/10})$, hence subexponential growth
or decay in depth.  Let $B_N^*$ denote the analogous product
for an inclined run.  The two mirror inclinations have identical
step coefficients and identical $B_N^*$ by the reflection
covariance of the bulk prescription.

\begin{proposition}[Comparison of normalization products]
\label{prop:source-products}
For the same bare coupling $\beta_N$ in the standard run and
either inclined run, there is a constant $C_{L,H}<\infty$,
independent of depth, such that
\begin{equation}
 C_{L,H}^{-1}\le\frac{B_N^*}{B_N}\le C_{L,H}.
 \label{eq:orig-33}
\end{equation}
\end{proposition}

\begin{proof}
Compare the bulk histories after their first steps.  Their
precise couplings satisfy $|\lambda_j|\le C_L$ by
Equation~\eqref{eq:orig-13}.  Put $x=N-1-j$.  The history
forcing in Equation~\eqref{eq:orig-5} is bounded by
$C_LH_N^D(A_0/L^2)^x$, after enlarging a fixed exponent $D$.
The coupling forcing is at most
$C_L(\log H_j)^D/H_j$.  Iteration of the shape inequality,
together with its fixed cap, gives
\begin{equation}
 u_j\le
 \min\{C,C_LH_N^D\rho_1^x\}
       +C_L\frac{(\log H_j)^D}{H_j}
 \label{eq:source-product-shape}
\end{equation}
for some $\rho_1<1$.  A slightly larger geometric base absorbs
geometric convolution factors.  The second forcing has this
same form after convolution: along the iteration the scale
$H_j$ decreases, and its logarithmic ratio is eventually
monotone, with finitely many initial values absorbed into the
constant.

Apply Equation~\eqref{eq:orig-29} and clip $u_s$ by its cap
in the same way.  Enlarging $C_L,D$ and $\rho_1<1$ if needed,
Equation~\eqref{eq:source-product-shape} also bounds $u_s$
and $|c_{j,N}^*-c_{j,N}|/R_j$.
The contribution from the second term is summable uniformly:
\[
 \sum_{j\ge0}R_j\frac{(\log H_j)^D}{H_j}
 =\sum_{j\ge0}\frac{(\log H_j)^D}{H_j^{11/10}}
 <\infty.
\]
For the clipped first term, split the iteration after
$J_N=\lceil K\log H_N\rceil$ steps, with $K$ fixed and large.
On the first $J_N$ steps, $H_j$ is comparable with $H_N$
for all sufficiently large $N$, so their total contribution
is at most $C_{L,H}H_N^{-1/10}\log H_N$.  Thereafter use
the unclipped geometric bound.  Since $R_j\le H^{-1/10}$,
this remaining sum is bounded by
\[
 C_{L,H}H_N^D\sum_{x\ge J_N}\rho_1^x
 \le C_{L,H}H_N^{D+K\log\rho_1}.
\]
Choose $K$ so that the last exponent is negative.  Finitely
many small depths cause no problem.  The single initial step
is controlled directly by Equation~\eqref{eq:orig-28}.
We have proved
\[
 \sum_{j=1}^N|c_{j,N}^*-c_{j,N}|\le C_{L,H}.
\]
All coefficients lie in a fixed positive neighborhood of one,
so the same estimate holds for the sum of their logarithmic
differences.  Exponentiating proves
Equation~\eqref{eq:orig-33}.  The proof does not require the
inclined terminal coupling to equal the terminal coupling
of the standard run.
\end{proof}
